% !PRINTED none
% Verlagsanzeige von Friedrich Vieweg & Sohn in Braunschweig; ohne Seitenzahl. Auf den
% Verlagskopf und einen kurzen Mittelstrich folgen sechs vollständige, durch kurze
% Mittelstriche getrennte Anzeigen:
% 1. „Was sind und was sollen die Zahlen?“ von Richard Dedekind, Professor der Mathematik an
% der technischen Hochschule zu Braunschweig; zweite unveränderte Auflage; 8.; geh.;
% Preis 1 Mark 60 Pfennig.
% 2. „Vorlesungen über Zahlentheorie“ von P. G. Lejeune-Dirichlet; herausgegeben und mit
% Zusätzen versehen von R. Dedekind, Professor an der technischen Hochschule
% Carolo-Wilhelmina zu Braunschweig; vierte umgearbeitete und vermehrte Auflage; gr. 8.;
% geh.; Preis 14 Mark.
% 3. „Lehrbuch der politischen Arithmetik“, für höhere Handelsschulen (Handelsakademien)
% und zum Selbstunterricht bearbeitet von F. S. Holzinger, Professor an der öffentlichen
% Handelsakademie in Linz; gr. 8.; geh.; Preis 3 Mark 50 Pfennig.
% 4. „Anleitung zur Durchmusterung des Himmels. Astronomische Objecte für gewöhnliche
% Teleskope. Ein Hand- und Hülfsbuch für alle Freunde der Himmelskunde, besonders für die
% Besitzer von Fernrohren“ von Dr. Hermann J. Klein; zweite verbesserte Auflage; mit 75
% Holzstichen, 5 Tafeln, zum Theil in Farbendruck, 4 Sternkarten und 1 Titelbilde; 8.; geh.;
% Preis 24 Mark.
% 5. „Sammlung von Formeln der reinen und angewandten Mathematik“ von Dr. W. Láska; mit
% drei Tafeln; gr. 8.; geh. 26 Mark; in Halbfranz geb. 28 Mark.
% 6. „Die Beobachtung der Sterne sonst und jetzt“ von J. Norman Lockyer, Mitglied der Royal
% Society, corr. Mitglied des Instituts von Frankreich; autorisirte deutsche Ausgabe;
% Uebersetzt von G. Siebert; mit 217 Holzstichen; 8.; geh.; Preis 18 Mark.
